\documentclass[11pt]{article}
\usepackage[margin=2.3cm]{geometry}
\usepackage{booktabs}
\usepackage{parskip}
\setlength{\parskip}{6pt}
\newcommand{\code}[1]{\texttt{\small #1}}
\title{\vspace{-1.2cm}\textbf{The tutor's four numbers}\\[3pt]
\large What they do, and what happens when they are wrong}
\author{}
\date{\small 28 September 2026}
\begin{document}
\maketitle
\vspace{-1.0cm}
\noindent\rule{\textwidth}{0.4pt}

\section*{1. What this part is about}

The tutor in our simulator does not watch the student learn. It only sees
answers, right or wrong, and from those it keeps a running guess about
whether each concept is known. That guess is built from four numbers.

This document is about those four numbers. It asks three things in order.
What does each number actually do? What goes wrong when the tutor's
numbers do not match the student's? And if students differ from each
other, what single set of numbers should the tutor use for everybody?

Everything here runs on a small practice curriculum of eight concepts
arranged in four layers, and every figure comes from a file of results
that one command regenerates.

\section*{2. The four numbers, in plain words}

% GENERATED by make_tables.py from sanity_outputs/. Do not edit:
% edit the experiment, rerun make_sanity.py, then make_tables.py.
\begin{center}
\small
\begin{tabular}{@{}cll@{}}
\toprule
 & \textbf{Meaning} & \textbf{Default}\\
\midrule
$p(L_0)$ & probability the concept is known before any question & 0.15\\
$p(T)$ & probability an unknown concept becomes known at an attempt & 0.20\\
$p(G)$ & probability of a correct answer while it is unknown & 0.25\\
$p(S)$ & probability of a wrong answer while it is known & 0.10\\
\bottomrule
\end{tabular}
\end{center}


\noindent Read them as sentences about a student:

\begin{itemize}\itemsep2pt
\item $p(L_0)$: some students already know a thing before you teach it.
\item $p(T)$: each time you ask about something they do not know, there
is a chance they pick it up.
\item $p(G)$: someone who does not know it can still answer correctly by
luck.
\item $p(S)$: someone who does know it can still slip and get it wrong.
\end{itemize}

\noindent The last two are why this is hard. A right answer is not proof
of knowing, and a wrong answer is not proof of not knowing. The tutor has
to weigh evidence rather than read it off.

Two things matter to a teacher, so those are what we measure. How long
does the tutor take to notice that a concept has been learned, counted in
questions asked about that concept. And how often does it announce
mastery too early, before the student has actually learned it.

\section*{3. What each number does}

Here the student and the tutor use the \emph{same} values, so the tutor is
not wrong about anything. We just change the numbers and watch.

% GENERATED by make_tables.py from sanity_outputs/. Do not edit:
% edit the experiment, rerun make_sanity.py, then make_tables.py.
\begin{center}
\small
\begin{tabular}{@{}lccc@{}}
\toprule
 & \textbf{swept over} & \textbf{questions to notice} & \textbf{declared too early}\\
\midrule
$p(L_0)$ & 0.05 to 0.50 & 3.1 to 2.4 & 0.050 to 0.075\\
$p(T)$ & 0.05 to 0.40 & 4.3 to 1.8 & 0.076 to 0.042\\
$p(G)$ & 0.05 to 0.45 & 2.1 to 3.8 & 0.000 to 0.058\\
$p(S)$ & 0.02 to 0.30 & 2.9 to 3.7 & 0.025 to 0.071\\
\bottomrule
\end{tabular}
\end{center}


\noindent Three of these behave the way common sense says they should.

When learning is faster, $p(T)$ large, the tutor notices sooner and
announces too early less often. Everything gets easier at once.

When guessing is easier, $p(G)$ large, lucky right answers look like
knowledge, so the tutor is fooled more often and takes longer to be sure.

When slipping is more likely, $p(S)$ large, a student who does know
something keeps answering wrongly, so the tutor holds back and takes
longer.

The fourth, $p(L_0)$, is the odd one. Raising it makes the tutor notice
faster but also announce too early more often. That is not a contradiction.
If many students already know things, starting optimistic is usually right
and occasionally badly wrong.

There is one catch worth stating. The tutor announces mastery when its
guess passes a bar, and passing that bar takes a whole number of correct
answers in a row. So these curves are steps, not smooth slopes. A change
in a parameter that does not move that whole number changes nothing at
all, and a change that does move it shifts everything at once. That is why
some of these lines bend the wrong way in places.

\section*{4. When the tutor's numbers are wrong: one student}

Now we keep the student fixed and give the tutor the wrong value on
purpose, by a known amount, one number at a time. If the tutor's guess is
honest, its error should be zero when it is right and grow as it is made
more wrong.

% GENERATED by make_tables.py from sanity_outputs/. Do not edit:
% edit the experiment, rerun make_sanity.py, then make_tables.py.
\begin{center}
\small
\begin{tabular}{@{}lcccc@{}}
\toprule
\textbf{tutor wrong about} & \textbf{gap when right} & \textbf{worst value} & \textbf{gap there} & \textbf{declared too early}\\
\midrule
$p(T)$ & +0.002 & 0.05 (-0.15) & +0.133 & 0.013\\
$p(G)$ & +0.002 & 0.05 (-0.20) & -0.063 & 0.171\\
$p(S)$ & +0.002 & 0.30 (+0.20) & -0.037 & 0.070\\
\bottomrule
\end{tabular}
\end{center}


\noindent That is what happens. The middle column is near zero for all
three, which is the check that the machinery is sound.

Beyond that, the three numbers do not matter equally. Being wrong about
$p(T)$ moves the tutor's belief about twice as much as being wrong about
either of the others. If you only have time to measure one thing about a
new group of students, measure how fast they learn.

The more useful finding is about direction. Being wrong about $p(G)$ and
being wrong about $p(T)$ push announcing-too-early in \emph{opposite}
ways. A tutor that thinks students guess more than they do becomes
cautious and announces less. A tutor that thinks students learn faster
than they do becomes hasty and announces more. So two different mistakes
can cancel, and a tutor can look well behaved while being wrong about both.

We deliberately do not rank $p(T)$ against $p(G)$ here. We tried, and the
ranking flipped depending on which side of zero we looked at and which
random seed we used. It is not a stable result and we are not reporting
one.

\section*{5. When every student is different}

Real students are not identical. Here each simulated student draws their
own values, and the tutor uses one set for everybody, sitting at the
middle of the group.

% GENERATED by make_tables.py from sanity_outputs/. Do not edit:
% edit the experiment, rerun make_sanity.py, then make_tables.py.
\begin{center}
\small
\begin{tabular}{@{}lccc@{}}
\toprule
\textbf{how different the students are} & \textbf{belief gap} & \textbf{declared too early} & \textbf{questions to notice}\\
\midrule
$\pm 0.00$ & +0.002 & 0.046 & 3.06\\
$\pm 0.05$ & +0.001 & 0.041 & 3.07\\
$\pm 0.10$ & -0.008 & 0.058 & 3.11\\
$\pm 0.15$ & -0.017 & 0.075 & 3.13\\
\bottomrule
\end{tabular}
\end{center}


\noindent The first row is the control: everybody matches the tutor.
Reading down, the tutor gets steadily worse, but slowly. Even when
students vary quite a lot, the tutor's belief is off by under two points
and it announces too early on under eight percent of concepts. One model
for a mixed group is workable, not ideal.

\section*{6. So what should the tutor assume?}

The obvious follow-up: if we cannot give every student their own numbers,
which single set is best, and how much are we losing by not adapting?

The second half needs something to compare against. So we also run an
\emph{oracle} tutor that is simply told each student's true values. No
real tutor can do this. It is there to mark the ceiling: whatever the
oracle gains is the most that adapting could ever be worth.

% GENERATED by make_tables.py from sanity_outputs/. Do not edit:
% edit the experiment, rerun make_sanity.py, then make_tables.py.
\begin{center}
\small
\begin{tabular}{@{}llcccccc@{}}
\toprule
\textbf{vary} & \textbf{spread} & \textbf{best fixed} & \textbf{its Brier} & \textbf{oracle} & \textbf{buys} & \textbf{95\% CI} & \textbf{resolved}\\
\midrule
p$\_$S & $\pm 0.05$ & 0.10 & 0.0835 & 0.0834 & $+0.0002$ & $\pm 0.0005$ & no\\
p$\_$S & $\pm 0.10$ & 0.10 & 0.0836 & 0.0834 & $+0.0002$ & $\pm 0.0013$ & no\\
p$\_$S & $\pm 0.15$ & 0.10 & 0.0854 & 0.0835 & $+0.0019$ & $\pm 0.0013$ & yes\\
p$\_$G & $\pm 0.05$ & 0.25 & 0.0826 & 0.0832 & $-0.0006$ & $\pm 0.0016$ & no\\
p$\_$G & $\pm 0.10$ & 0.25 & 0.0824 & 0.0835 & $-0.0011$ & $\pm 0.0018$ & no\\
p$\_$G & $\pm 0.15$ & 0.25 & 0.0829 & 0.0828 & $+0.0001$ & $\pm 0.0020$ & no\\
\bottomrule
\end{tabular}
\end{center}


\noindent Two answers.

The best fixed tutor is always the one sitting at the middle of the group.
We could not find a value that beat it, at any spread, for either number.
So there is no clever offset to discover.

And adapting is worth very little. Knowing each student's own $p(S)$ helps
only when students vary a lot, and then only by about two percent. Knowing
each student's own $p(G)$ never helped measurably at all.

One honest caveat. ``No measurable difference'' is not the same as ``no
difference''. We did not set a threshold in advance for how small a gain
would count as none. All we can say is that if adapting helps, it helps by
less than this study could see.

We should also say why we score this on the Brier score rather than
something simpler. We first used the average gap between what the tutor
believes and what is true, and that was a mistake. A tutor that is too
confident about half its students and not confident enough about the other
half scores perfectly on that measure, because the two halves cancel. That
is exactly the error adapting would fix, so the measure was blind to the
question being asked. The Brier score does not cancel, and it correctly
picks out the true value when we feed it one.

\section*{7. What we are not claiming}

All of this is on one small curriculum of eight concepts. Nothing has been
rerun on the larger one, so none of it is established beyond this graph.

Every number describes a simulated student, not a person. No result here
has been checked against a real learner.

The bar for announcing mastery is fixed at 0.9 and is not swept, and since
it is a whole-number-of-correct-answers effect, some of these comparisons
would move if the bar moved.

\section*{8. Running it again}

\code{python3 make\_sanity.py} reruns every study and writes the result
files. \code{python3 make\_tables.py} turns those files into every table in
this document, so no number here is typed by hand.
The command \code{python3 make\_tables.py -{}-check} fails if the document
and the results ever drift apart. Seeds are fixed, so a rerun gives the
same numbers.

\end{document}
